\chapter{Simulation--Production Parity}\label{ch:pl:simulation-production-parity}

A strategy that had passed its replay parity test quoted a tick wider than intended on its first morning in production, for the whole
morning. Its opening period ended on a timer, and the timer had been set from the machine's clock: in production the machine's clock
and the venue's agreed, in the backtest they did not, and the parity test had compared the backtest with itself. The code was the same
object in both places; the environments around it were not. This chapter builds the host that lets one strategy object run in a
backtest, in the simulator and on a production-shaped path, and the harness that finds, event by event, where two of those runs part.

\section{One strategy, several environments}

\begin{definition}[Simulation--production parity]\label{def:pl:simulation-production-parity:parity}
\emph{Simulation--production parity}\index{simulation--production parity} is the property that a strategy, given the same inputs in
the same order, produces the same outputs in simulation and in production: the same orders and cancels, at the same prices and
quantities, in response to the same events.
\end{definition}

Parity is a property of the strategy and of everything around it. Book~7 (chapter~19) tested it by replaying a live day's inputs into
the backtest and comparing orders; Book~12 (chapter~24) met its cousin in machine learning, the training--serving skew. A strategy can
fail it by itself (it reads something that is not an input: the machine's clock, a random number, the order of a dictionary) or through
its environments (they deliver the same market in different types, orders or semantics). The platform removes the second kind by
construction and exposes the first kind by test.

\begin{definition}[Strategy host, environment adapter]\label{def:pl:simulation-production-parity:host}
A \emph{strategy host}\index{strategy host} is the component that holds a strategy object, gives it its only view of the world --- a
clock, market data and order entry --- and delivers its inputs in the platform's \omterm{def:pl:backtest-engine-architecture:event}{event model}. An \emph{environment
adapter}\index{environment adapter} connects the host to one environment (a backtest, a simulator, a production connection), translating
that environment's events into the host's inputs and the strategy's requests into the environment's messages.
\end{definition}

\begin{omfigure}
\begin{tikzpicture}[font=\scriptsize,box/.style={draw,rounded corners=2pt,align=center,minimum height=0.7cm,inner sep=3pt},>=Stealth]
\node[box,fill=omMeth!15,minimum width=3.4cm] (s) at (0,2.5) {strategy object\\ (the same in every environment)};
\node[box,fill=omDef!15,minimum width=10.4cm] (h) at (0,1.2) {strategy host: clock $\cdot$ market data $\cdot$ order entry $\cdot$
inputs delivered in the event model};
\node[box,fill=omThm!10,minimum width=3.1cm] (a1) at (-4.0,-0.4) {simulator in process\\ agent context, objects};
\node[box,fill=omThm!10,minimum width=3.1cm] (a2) at (0,-0.4) {production-shaped\\ gateway, bytes, journal};
\node[box,fill=omThm!10,minimum width=3.1cm] (a3) at (4.0,-0.4) {replay\\ a production journal};
\draw[<->] (s) -- (h);
\foreach \a in {a1,a2,a3} {\draw[<->] (\a.north) -- (\a.north|-h.south);}
\node[box,fill=gray!10,minimum width=5.2cm] (x) at (-1.8,-1.9) {\texttt{firm.exchsim}: one venue, the same session};
\draw[->] (a1.south) -- (a1.south|-x.north);
\draw[->] (a2.south) -- (a2.south|-x.north);
\draw[->,dashed] (a2.east) -- node[above,font=\tiny]{journal} (a3.west);
\end{tikzpicture}

\omcaption{One strategy object in three environments. The host gives it a clock, market data and order entry, and delivers every input
in the \omterm{def:pl:backtest-engine-architecture:event}{event model}; each \omterm{def:pl:simulation-production-parity:host}{environment adapter} translates its environment. The in-process and production-shaped adapters trade on the same
simulated session; the replay reads the production run's input journal.}\label{fig:pl:simulation-production-parity:host}
\end{omfigure}

The chapter's three environments trade on Book~10's simulator with one \texttt{firm.tape} hour as background flow (chapter~11). The
first calls the simulator's agent context directly. The second is shaped like production: the strategy's orders pass through Book~13's
order gateway (\texttt{firm.ordergw}: its order state machine, its throttle and its worst-case exposure check), are encoded as the
simulator's order-entry messages in SoupBinTCP frames on an in-memory byte stream, and the venue's reports come back as bytes the same way;
every input the strategy receives is written to an input journal in Book~13's format (chapter~23: a receive time and a 48-byte record). The
third replays that journal into a fresh strategy object: the backtest of recorded inputs.

\section{Abstracting time}

The host's clock is the only clock a strategy may read: \texttt{ctx.now} is the time of the input being delivered, and
\texttt{ctx.set\_timer(delay, tag)} asks for a timer input \texttt{delay} later. In the simulator both are simulated time; in production
both would be the venue-synchronised clock of the machine; in the replay, \texttt{now} is the journal's receive time and timers are not set
but replayed, since the timers that fired in production are inputs in the journal. A strategy that reads the machine's clock directly reads
something the host does not control, and its behaviour then depends on where and how fast it runs. The chapter's harness makes this
visible by giving each environment the machine clock it would really have: the venue's time in production, and, in the simulator and the
replay, the clock of a machine replaying the session at 20:00, a hundred times faster than real time.

\section{Abstracting market data}

\begin{definition}[Market-data abstraction]\label{def:pl:simulation-production-parity:md}
A \emph{market-data abstraction}\index{market-data abstraction} is the one representation of market events that a hosted strategy
receives in every environment: here the top of the book as integers --- bid, bid size, ask, ask size --- with prices in 1/10\,000 of the
currency unit, whatever the source's own encoding.
\end{definition}

The rule that prices are integers is Book~10's and the series': a price is a count of the smallest unit, and every conversion to a binary
floating-point number is a chance to be off by one. Python's documentation says it plainly: most decimal fractions cannot be represented
exactly as binary fractions. \texttt{99.99 - 0.01} is \texttt{99.97999999999999}, and \texttt{int()} of ten thousand times it is
999\,799, one unit below the tick grid of 999\,800.

\section{Abstracting order entry}

\begin{definition}[Order-entry abstraction]\label{def:pl:simulation-production-parity:oe}
An \emph{order-entry abstraction}\index{order-entry abstraction} is the one interface through which a hosted strategy sends and cancels
orders and learns their fate, with the same semantics in every environment: which orders count as working (here, sent and not yet known to
be finished, whether or not acknowledged), which callbacks arrive (acknowledgement, fill, cancel, reject), and what happens to a request that
can no longer succeed.
\end{definition}

Semantics are where environments differ silently. Book~13's gateway (chapter~21) refuses to cancel an order it already knows to be filled;
a simulator's agent context passes the cancel to the venue, which rejects it as too late. The same strategy then sees a reject in one
environment and nothing in the other, and the venue sees an extra message that takes 500~nanoseconds of its matching engine's time. The
harness below found exactly this difference on its first clean run, before any defect had been planted: the in-process adapter now refuses,
as the gateway does, to cancel an order a report has already finished.

\omcode{code/firm/strathost/firm_strathost.py}{125}{145}{The host's side of the abstractions: which orders are working, timers, and
the delivery of simultaneous inputs in the event model (market data, then order entry, then timers).}\label{lst:pl:simulation-production-parity:host}

\omcode{code/firm/strathost/firm_strathost.py}{278}{305}{The production-shaped adapter: an order passes Book~13's gateway, is encoded
and framed onto a byte stream, and becomes an order at the venue side of the stream.}\label{lst:pl:simulation-production-parity:prod}

\section{The parity harness}

\begin{definition}[Parity harness]\label{def:pl:simulation-production-parity:harness}
A \emph{parity harness}\index{parity harness} runs one strategy in two environments on the same inputs and compares their outputs
event by event, reporting the first output that differs, the inputs delivered before it, and how many outputs differ in all.
\end{definition}

The harness compares the production-shaped run with the in-process simulator run (the same session, two adapters) and with the replay of
its own journal (the same inputs, delivered again). Outputs are compared as the venue would see them: time, order or cancel, side, price on
the wire, quantity. The first difference and the number of inputs before it are what a developer needs to start reading; the count of
differing outputs says how much a defect costs in an hour.

\omcode{code/firm/strathost/firm_strathost.py}{431}{448}{The comparison: the first differing output, the inputs delivered before it,
and the outputs of one run with no identical output in the other.}\label{lst:pl:simulation-production-parity:compare}

The chapter's strategy (\cref{lst:pl:simulation-production-parity:quoter}) quotes a lot at the touch, one tick wider during an opening
period of thirty seconds from the start of the session, which a timer ends; it stops adding at three lots of inventory. Four defects are
planted, one at a time, in the production-shaped run:
\begin{itemize}
\item \emph{wall clock}: the strategy takes the end of its opening period and its timer's delay from the machine's clock;
\item \emph{float prices}: the production adapters hand prices to the strategy as floating-point currency and convert them back with
\texttt{int()};
\item \emph{callback order}: the production host delivers simultaneous inputs as its adapter read them, order entry first;
\item \emph{ack semantics}: production's \texttt{working()} lists only acknowledged orders.
\end{itemize}

\omcode{code/platforms/12-simulation-production-parity/python/pl_parity.py}{26}{60}{The hosted strategy. With the wall-clock defect on,
\texttt{\_now} reads the machine and the timer's delay mixes the machine's clock with the host's.}\label{lst:pl:simulation-production-parity:quoter}

\section{Four defects}

\begin{omfigure}
\begin{tikzpicture}
\begin{axis}[width=11cm,height=5.6cm,xbar,bar width=7pt,xmode=log,xmin=1,xmax=40000,log origin=infty,
  ytick=data,symbolic y coords={ack semantics,callback order,float prices,wall clock},enlarge y limits=0.2,
  xlabel={count (log scale)},tick label style={font=\small},label style={font=\small},grid=major,grid style={gray!20},
  legend style={font=\footnotesize,at={(0.5,1.03)},anchor=south,legend columns=2},table/col sep=comma,
  nodes near coords,point meta=explicit symbolic,nodes near coords style={font=\scriptsize,anchor=west}]
\addplot[fill=omDef!45,draw=omDef] table[x=events,y=defect,meta=events]{figdata/platforms/12-simulation-production-parity/defects_sim.csv};
\addlegendentry{inputs before the first diverging output}
\addplot[fill=omThm!45,draw=omThm] table[x=affected,y=defect,meta=affected]{figdata/platforms/12-simulation-production-parity/defects_sim.csv};
\addlegendentry{outputs affected in the hour}
\end{axis}
\end{tikzpicture}

\omcaption{Each planted defect in the production-shaped run against the in-process simulator, on the same one-hour session: the inputs
delivered before the first output that differs, and the outputs (orders and cancels) of the production run with no identical output in the
simulator run. Data: \texttt{fig\_parity.py}.}\label{fig:pl:simulation-production-parity:defects}
\end{omfigure}

\Cref{fig:pl:simulation-production-parity:defects} and \cref{tab:pl:simulation-production-parity:defects} give the harness's report
for each defect on seed~1.

\begin{table}[ht]
\centering\small
\begin{tabular}{lrrrrr}
\toprule
defect & \multicolumn{2}{c}{first diverging output} & inputs before & \multicolumn{2}{c}{outputs affected} \\
 & index & seconds & & vs simulator & vs replay \\
\midrule
wall clock & 14 & 30.0 & 175 & 539 of 553 & 305 \\
float prices & 3 & 24.4 & 126 & 1\,686 of 1\,690 & 1\,161 \\
callback order & 20 & 33.0 & 203 & 387 of 463 & 416 \\
ack semantics & 2 & 1.0 & 2 & 8\,449 of 8\,452 & 8\,450 \\
\bottomrule
\end{tabular}
\caption{The harness's report for each planted defect: the index and time (seconds from the start of the session, one second before the
open) of the first output of the production-shaped run that differs from the in-process simulator's, the inputs delivered before it, and the
outputs of the production run with no identical counterpart in the simulator run and in the replay of its own journal.}
\label{tab:pl:simulation-production-parity:defects}
\end{table}

\begin{example}[The morning the timer fired]\label{ex:pl:simulation-production-parity:timer}
With the wall-clock defect, production is unaffected: its machine clock is the venue's, the timer fires thirty seconds after the start and
the strategy moves to the touch. In the simulator the timer's delay is the machine's 20:00 plus thirty seconds minus the session's 09:29:59,
more than ten hours, and it never fires in the hour; the machine clock, a hundred times slower than the session's, keeps the opening period
open for fifty minutes. The first diverging output is production's cancel at 30.0~seconds, after 175 inputs, and 539 of production's 553
outputs have no twin in the simulator. The replay of production's own journal receives the timer from the journal but still reads its own
machine clock, requotes a tick wide, and parts at 30.1~seconds.
\end{example}

\begin{omfigure}
\begin{tikzpicture}[font=\scriptsize,>=Stealth,x=0.19cm]
\foreach \y/\l in {0/production,-1.1/simulator,-2.2/replay} {\node[anchor=east] at (-1,\y) {\l};
  \draw[gray] (0,\y) -- (60,\y);}
\draw[->] (0,-3.0) -- (61.5,-3.0) node[right]{minutes};
\foreach \m in {0,10,20,30,40,50,60} {\draw (\m,-2.9) -- (\m,-3.1) node[below]{\m};}
\fill[omThm!50] (0,-0.15) rectangle (0.5,0.15);
\draw[omThm,thick] (0.5,0.35) -- (0.5,-0.15) node[pos=0,above,font=\tiny]{timer: 30 s};
\fill[omDef!40] (0.5,-0.12) rectangle (60,0.12);
\fill[omThm!50] (0,-1.25) rectangle (50,-0.95);
\fill[omDef!40] (50,-1.22) rectangle (60,-0.98);
\node[font=\tiny] at (25,-0.75) {a tick wide for 50 minutes; the timer (10.5 hours) never fires};
\fill[omThm!50] (0,-2.35) rectangle (50,-2.05);
\fill[omDef!40] (50,-2.32) rectangle (60,-2.08);
\draw[omThm,thick] (0.5,-1.85) -- (0.5,-2.35);
\node[font=\tiny] at (25,-1.85) {the journalled timer fires; the machine clock keeps it wide};
\end{tikzpicture}

\omcaption{The wall-clock defect in the three environments over the hour: shaded dark, the strategy quotes a tick wide; light, at the
touch. In production the machine's clock is the venue's and the opening lasts thirty seconds. In the simulator and the replay the machine
runs a hundred times slower than the session, from 20:00, and the opening lasts fifty minutes.}\label{fig:pl:simulation-production-parity:clock}
\end{omfigure}

The other three defects read the same way. With float prices the first difference is the third output, at 24.4~seconds: the strategy's wide
bid $99.99 - 0.01$ leaves as 999\,799; 1\,161 of production's 1\,501 orders are off the tick grid, rejected by the venue and sent again at the
next quote. With inputs delivered in arrival order the strategy first differs at 33.0~seconds, after 203 inputs, on an instant where a fill
and a quote change arrived together: delivered fill first, it sent its next sell order at once; delivered as the \omterm{def:pl:backtest-engine-architecture:event}{event model} delivers them,
quote first, it sent the same order 0.19~seconds later. Neither order is wrong; two orders are. With working meaning acknowledged, the next
quote change after the first two orders, half a microsecond later and long before their acknowledgements, makes the strategy send both
again, and every quote change inside the forty-microsecond round trip does the same: 8\,452 outputs in the hour, 8\,406 of them refused by
the gateway's worst-case exposure limit, which is all that stops the storm.

\begin{remark}[Why the replay alone was not enough]\label{rem:pl:simulation-production-parity:replay}
A replay parity test (Book~7, chapter~19) catches a strategy that reads what is not an input, but only if the replay environment differs from
production in that input: it caught the wall clock here because the harness gave the replay a machine clock of its own. It catches adapter
defects only when the replay host has the correct semantics and production's does not. Comparing two live-shaped environments on the same
session catches what the replay misses; comparing the replay with production catches what two simulators share.
\end{remark}

\section{Parity after the fixes}

With every defect fixed --- the strategy reads only the host's clock, prices are integers end to end, the host delivers simultaneous inputs
in the \omterm{def:pl:backtest-engine-architecture:event}{event model} everywhere, and working means sent and not finished everywhere --- the three environments agree on every output for ten
one-hour sessions (\cref{tab:pl:simulation-production-parity:clean}): 8\,887 outputs from 188\,079 inputs, no difference between the
production-shaped run, the simulator and the replay.

\begin{table}[ht]
\centering\small
\begin{tabular}{rrrcc}
\toprule
seed & inputs & outputs & simulator & replay \\
\midrule
1 & 9\,492 & 553 & identical & identical \\
2 & 16\,085 & 832 & identical & identical \\
3 & 36\,396 & 1\,463 & identical & identical \\
4 & 37\,621 & 1\,359 & identical & identical \\
5 & 13\,911 & 784 & identical & identical \\
6 & 29\,785 & 1\,426 & identical & identical \\
7 & 6\,947 & 409 & identical & identical \\
8 & 13\,410 & 726 & identical & identical \\
9 & 20\,027 & 979 & identical & identical \\
10 & 4\,405 & 356 & identical & identical \\
\bottomrule
\end{tabular}
\caption{Parity after the fixes: for ten one-hour sessions, the inputs and outputs of the production-shaped run and whether the in-process
simulator run and the replay of the production journal produced identical outputs.}\label{tab:pl:simulation-production-parity:clean}
\end{table}

\section{Tutorial: one strategy, three environments, four defects}\label{tut:pl:simulation-production-parity:tut}

\begin{tutorial}
\textbf{Goal.} Host one strategy in three environments, plant four parity defects and let the harness find each; then show parity.
\textbf{End state:} \cref{tab:pl:simulation-production-parity:defects,tab:pl:simulation-production-parity:clean}.
\begin{tutsteps}
\item \textbf{Environments}: \texttt{run\_env} with \texttt{prod} and with \texttt{sim} on the same session; \texttt{replay} of the
production run's journal.
\item \textbf{The harness}: \texttt{compare(prod, sim)} and \texttt{compare(prod, replay)}.
\item \textbf{Defects}: \texttt{pl\_parity.defect\_table()}, one defect at a time.
\item \textbf{Parity}: \texttt{pl\_parity.clean()} on seeds 1 to 10.
\end{tutsteps}
\textbf{What to change next.} Plant a fifth defect --- a strategy that iterates over a set of order identifiers when it cancels --- and see
whether the harness finds it; give the in-process adapter a different latency from production's and see what parity then means.
\end{tutorial}

\section{Build: the strategy host}\label{bld:pl:simulation-production-parity:strathost}

\begin{build}
\bfield{Purpose} One strategy object that runs unchanged in the backtest, the simulator and production, and a harness that proves it.
\bfield{Interface} \texttt{HostedStrategy} (\texttt{on\_start}, \texttt{on\_book}, \texttt{on\_ack}, \texttt{on\_fill},
\texttt{on\_cancel}, \texttt{on\_reject}, \texttt{on\_timer}); the context (\texttt{now}, \texttt{set\_timer}, \texttt{send},
\texttt{cancel}, \texttt{working}, \texttt{position}); \texttt{run\_env(env, strategy, seconds, seed, defects)};
\texttt{replay(journal, strategy)}; \texttt{compare(a, b) -> Parity}; \texttt{Defects}.
\bfield{Rules} The strategy reads nothing the host does not give it; prices are integers; simultaneous inputs in the \omterm{def:pl:backtest-engine-architecture:event}{event model}; the same
order-entry semantics in every adapter; production's inputs journalled in Book~13's format; Book~13's gateway and Book~10's simulator used,
not edited.
\bfield{Acceptance tests} \texttt{code/firm/strathost/tests/}: the delivery order of simultaneous inputs, with and without the defect; clean
parity of production, simulator and replay on a short session, and the journal's records; divergence under each adapter defect; the machine
clock; the float adapter's truncation; no cancel of a finished order.
\bfield{Stretch} A fourth environment on a real socket; parity across implementation languages (the C++20 and Rust gateways of Book~13); a
bisecting harness on Book~13's \texttt{first\_difference}.
\end{build}

\begin{omsources}
\item One Quant Book~7, chapter~19 (replay parity); One Quant Book~13, chapters 21 and 23 (the order gateway, input journals); One Quant
Book~10, chapter~26 (the simulator and its protocols).
\item Python documentation, ``Floating-Point Arithmetic: Issues and Limitations''.
\end{omsources}

\section{Exercises}

\begin{exercise}[$\star$]\label{exo:pl:simulation-production-parity:1}
Why is \texttt{int(ten thousand times (99.99 - 0.01))} 999\,799, and what is the correct conversion?
\end{exercise}

\begin{exercise}[$\star$]\label{exo:pl:simulation-production-parity:2}
In the replay environment, why are timers not set but replayed?
\end{exercise}

\begin{exercise}[$\star$]\label{exo:pl:simulation-production-parity:3}
With the wall-clock defect, why is production itself unaffected?
\end{exercise}

\begin{exercise}[$\star\star$]\label{exo:pl:simulation-production-parity:4}
Why does the ack-semantics defect produce a storm rather than a few duplicate orders, and what stops it?
\end{exercise}

\begin{exercise}[$\star\star$]\label{exo:pl:simulation-production-parity:5}
The harness found an unplanned difference on its first clean run. Describe it, and say which adapter was right.
\end{exercise}

\begin{exercise}[$\star\star$]\label{exo:pl:simulation-production-parity:6}
Why can the count of outputs affected differ between the comparison with the simulator and with the replay?
\end{exercise}

\begin{exercise}[$\star\star\star$]\label{exo:pl:simulation-production-parity:7}
\emph{Coding.} Run \texttt{run\_three} for seed~2 with each defect. Are the first diverging outputs at the same times as for seed~1?
\end{exercise}

\begin{exercise}[$\star\star\star$]\label{exo:pl:simulation-production-parity:8}
\emph{Find the flaw.} ``Our parity test replays yesterday's backtest inputs into the backtest and gets identical orders, so the strategy
is ready for production.''
\end{exercise}

\section{Problem: The Morning the Timer Fired}

\begin{problem}[{Weekend problem --- four defects and a clean bill}]\label{pb:pl:simulation-production-parity:1}
The strategy of \cref{lst:pl:simulation-production-parity:quoter}, three environments, one-hour sessions.

\medskip\noindent\textbf{Part I --- The host.}
\begin{enumerate}
\item What may a hosted strategy read, and through what?
\item Why does the host buffer inputs by instant?
\item What does the production-shaped adapter add to the in-process one?
\item What does the input journal hold, and in which format?
\item What does the replay environment do with the strategy's orders?
\end{enumerate}

\medskip\noindent\textbf{Part II --- The defects.}
\begin{enumerate}[resume]
\item How does the wall-clock defect show in the simulator and in the replay?
\item What does the float-price defect do to the wide bid?
\item When does the callback-order defect matter?
\item What does ack semantics change, and why the storm?
\item Which of the four would a replay of production's journal into a correct host miss, and why?
\end{enumerate}

\medskip\noindent\textbf{Part III --- The report.}
\begin{enumerate}[resume]
\item Give the inputs before the first diverging output for each defect.
\item Give the outputs affected in the hour for each.
\item Which defect is the fastest to show, and which the slowest?
\item What unplanned difference did the harness find?
\item How many outputs and inputs do the ten clean sessions have, and how many differences?
\end{enumerate}

\medskip\noindent\textbf{Part IV --- The verdict.}
\begin{enumerate}[resume]
\item State the \emph{named result}: the inputs until each defect's first diverging order, the orders affected in an hour, and the
parity after the fixes.
\item What would you require of every strategy before production?
\item What would you require of every \omterm{def:pl:simulation-production-parity:host}{environment adapter}?
\item Why compare two environments on the same session as well as replaying a journal?
\item In one sentence: what makes parity a property of the platform rather than of each strategy?
\end{enumerate}
\end{problem}

\section{Interview questions}

\begin{interviewq}[$\star$ \iqroles{developer}]\label{iq:pl:simulation-production-parity:1}
Your strategy behaves differently in production from the backtest on the same day. Where do you look first?
\end{interviewq}

\begin{interviewq}[$\star\star$ \iqroles{developer}]\label{iq:pl:simulation-production-parity:2}
Why should a trading system represent prices as integers?
\end{interviewq}

\begin{interviewq}[$\star\star$ \iqroles{developer, researcher}]\label{iq:pl:simulation-production-parity:3}
How would you make a strategy's use of time testable?
\end{interviewq}

\begin{interviewq}[$\star\star$ \iqroles{developer}]\label{iq:pl:simulation-production-parity:4}
A quote update and a fill have the same timestamp. Which should the strategy see first, and why does it matter?
\end{interviewq}

\begin{interviewq}[$\star\star\star$ \iqroles{developer}]\label{iq:pl:simulation-production-parity:5}
Design a harness that proves a strategy behaves identically in simulation and production.
\end{interviewq}
